\begin{table}[t]
\centering
\small
\caption{Sample-size sensitivity from 100 to 1000 examples per condition where available.}
\label{tab:sample-size}
\begin{threeparttable}
\begin{tabular}{llcccccc}
\toprule
n & Feature family & All mean & All SD & All runs & Competent mean & Competent SD & Competent runs \\
\midrule
100 & Confidence/option & 0.720 & 0.122 & 100 & 0.754 & 0.098 & 80 \\
100 & Conf./option + PCA hidden & 0.683 & 0.147 & 100 & 0.726 & 0.121 & 80 \\
300 & Confidence/option & 0.730 & 0.107 & 75 & 0.772 & 0.066 & 60 \\
300 & Conf./option + PCA hidden & 0.688 & 0.094 & 75 & 0.721 & 0.064 & 60 \\
500 & Confidence/option & 0.738 & 0.104 & 75 & 0.781 & 0.057 & 60 \\
500 & Conf./option + PCA hidden & 0.711 & 0.088 & 75 & 0.748 & 0.047 & 60 \\
1000 & Confidence/option & 0.745 & 0.093 & 75 & 0.783 & 0.051 & 60 \\
1000 & Conf./option + PCA hidden & 0.731 & 0.094 & 75 & 0.773 & 0.045 & 60 \\
\bottomrule
\end{tabular}

\begin{tablenotes}[flushleft]\footnotesize\item \parbox{0.95\textwidth}{ARC-Challenge has 299 validation examples; larger n values are available for the remaining datasets. Increasing sample size narrows variability but does not make hidden-augmented probes clearly dominate cheap features.}\end{tablenotes}
\end{threeparttable}
\end{table}
